\section{引言}

自主水下航行器（AUV）在海洋探索、水下检测和环境监测中发挥着日益重要的作用\cite{fossen2011handbook}。在轨迹跟踪任务中，海流是影响控制精度的主要扰动源。精确估计海流速度以用于前馈补偿，是实现鲁棒AUV控制性能的关键\cite{borhaug2007model}。

现有的AUV海流估计方法大致可分为两类。第一类为运动学海流观测器\cite{liang2018three,he2024three}，通过对位置或速度误差进行积分来估计惯性系海流分量。这类方法不依赖水动力模型精度，天然具有鲁棒性，但其收敛速度受限于位置误差的累积速率，通常需要数十秒，对突变海流响应缓慢。第二类为基于扩张状态观测器（ESO）的集总扰动估计\cite{xie2020extended}，将海流、模型不确定性和未建模动态统一估计为"总扰动"信号。虽然响应迅速，但ESO方法产生的是补偿信号而非物理可解释的海流估计，无法区分物理海流效应与模型误差。

两类方法面临的共同实际挑战是\textbf{水动力模型失配}。运动学观测器天然免疫，因其不使用动力学模型。ESO方法将模型误差无差别地吸收到集总扰动中。然而，当研究者尝试利用动力学模型进行直接物理海流估计时——如Kim等的高增益观测器配合在线递推最小二乘辨识\cite{kim2018estimating}，或Børhaug的非线性Luenberger观测器\cite{borhaug2007model}（假设完全已知动力学模型）——模型失配会直接污染海流估计。这引出了本文的核心问题：\textbf{能否设计一种既能利用动力学模型信息，又对模型失配具有鲁棒性的海流观测器？}

我们给出肯定的回答。我们的方法以XHY AUV具备CFD（RANS $k$-$\omega$ SST）预标定水动力模型（各自由度阻力系数拟合优度$R^2>0.99$，经静水拖曳试验和自由航行试验独立验证）的工程条件为基础。这提供的不仅是一个名义模型，而且是一个具有可量化精度边界的\textbf{确定性先验}。

在此基础上，本文提出CFD先验海流观测器（PC-RCO），采用预测-校正架构，包含两个关键创新：

\begin{enumerate}[itemsep=0pt,topsep=4pt]
    \item \textbf{连续Gauss-Markov（GM）正则化}：PC-RCO不依赖二值激励门控，而是在\textbf{每个}时间步施加弱GM衰减。高激励阶段梯度更新主导；低激励或高噪声阶段GM正则化逐渐将估计拉向先验均值，防止漂移并保持有界更新。这种连续混合消除了阈值门控的需求，实验证明后者无法在反馈控制AUV中可靠区分激励水平。
    \item \textbf{在线噪声估计的自适应更新率限幅}：PC-RCO根据速度预测残差维持指数滑动平均（EMA）的瞬时DVL噪声估计，动态调整单步更新上限（$\Delta c_{\max}$）——低噪声时紧限幅保证平滑收敛，高噪声时适度放宽以维持噪声穿透能力。绝对上限（0.20 m/s）确保任何条件下不会出现无界更新。
\end{enumerate}

我们在XHY AUV仿真平台上进行了系统的数值验证，涵盖6个CFD阻力模型失配水平（0--50\%）、4种基线方法（跨越3个方法论类别）和受控的DVL噪声条件。结果表明：（1）PC-RCO在30\% CFD阻力扰动下仅退化42\%，而EKF\cite{long2021trajectory}退化77\%，在零失配下精度比运动学基线高3倍（RMSE$_{V_c}=0.064$ vs.\ $0.192$）；（2）连续GM正则化在模型失配和高噪声条件下单独贡献57--218\%的精度改善，消融实验确认；（3）自适应限幅在高DVL噪声（$\sigma=0.10$ m/s）下将单步最大估计跳变降低34倍，防止物理不可能的瞬态估计（单步3 m/s跳变），而不牺牲稳态精度。

本文其余部分结构如下。第2节建立AUV动力学模型、海流参数化和CFD不确定性表征。第3节详述PC-RCO观测器设计。第4节给出系统的数值验证。第5节讨论适用性、局限性和部署考虑。第6节总结全文。
